
\chapter{Method}\label{ch:methodology}

This chapter specifies the method in three parts.
Section~\ref{sec:arch-requirements} states what the partition must
satisfy for the arena to mean anything, points to the measured
evidence that the induced partition satisfies it (reported with the
results, Section~\ref{sec:res-construction}), and names the three
things it does not guarantee; the construction formalism itself is in
Appendix~\ref{app:dag-formalism}.
Section~\ref{sec:arch-arena} specifies the arena --- what it judges,
what the judge sees, and how verdicts become rankings --- and points
to the inventory of mechanisms that have no implementation
(Appendix~\ref{app:not-implemented}).
Section~\ref{sec:meth-discipline} states the four measurement rules
this thesis holds itself to, each with the failure that produced it, so
that the results of Chapter~\ref{ch:results} are read against a
standard the reader already holds.

The ordering is deliberate. The partition is an instrument; the
question is what the arena measures with it. The construction receives
three pages here and eighteen in the appendix, in that proportion, for
the measured reason of Section~\ref{sec:res-granularity}: its tuning
has no measured evaluative consequence.

\section{What a Good Partition Needs}\label{sec:arch-requirements}

The arena judges responses inside cells and reports a ranking per cell.
For that to mean anything, four things must hold. They are stated here
as requirements, before any construction detail, so that the
construction can be assessed against them rather than described for its
own sake.

\begin{description}
\item[R1 --- Coherence.] Cells must be homogeneous enough that a single
rubric fits every query in the cell. This is the requirement the whole
design exists to serve (Section~\ref{sec:problem-statement}).
\item[R2 --- Support.] Cells must be large enough that per-cell judging
is identifiable --- the comparison graph connects --- and that per-cell
reliability is not dominated by sampling noise.
\item[R3 --- Stability.] The same corpus must yield the same cells, and
the construction must reach a state that one further application of its
own operator would not change.
\item[R4 --- Verifiability.] What the construction guarantees, and what
it does not, must be stated in advance of any result that depends on
it.
\end{description}

\subsection{How the Requirements Are Met}\label{sec:arch-taxonomy}

The construction is one paragraph in the main text.
The 14 MMLU-Pro domain labels seed 14 immutable top-level anchor nodes.
Every corpus query is embedded, and a fixed 256-dimensional prefix
slice of the embedding is $\ell_2$-normalised to the unit hypersphere
and used at every depth. Queries then migrate by embedding geometry, so
a query labelled Mathematics may come to sit under a different anchor;
this is what makes the induced (\textit{adapted}) partition differ from
the editorial (\textit{canonical}) one.
Each anchor is recursively split under a chance-corrected separation
score $J$, computed against random relabelling into the same size
profile. A split is accepted only when every child pair clears a
separation bar of $0.025$ \emph{and} the improvement in $J$ exceeds
twice its own paired bootstrap standard error. Two coarsening moves
--- absorption of below-floor children and merging of insufficiently
separated sibling pairs --- run in the same loop, so every operator maps
a tree to a tree. The loop runs until a certificate reports that one
more iteration would change nothing.
The formalism, the parameter derivations, the routing rule and the
certificate are given in full in Appendix~\ref{app:dag-formalism} and
Appendix~\ref{app:numerics}.

The resulting frozen artifact --- the single construction every result
in this thesis is attributed to --- is a tree of 154 nodes with 87
labelled, rubric-carrying leaves; its full statistics, the fixed-point
certificate, and the evidence per requirement are reported once, with
the results (Section~\ref{sec:res-construction}).

\subsection{Where the Evidence Is}\label{sec:arch-evidence-pointer}

Whether the construction meets R1--R4 is an empirical question, and
the evidence is reported with the results rather than here:
Table~\ref{tab:requirements-evidence}
(Section~\ref{sec:res-construction}) collects one line of evidence per
requirement, each marked \textsc{measured} or \textsc{argued} and each
with the caveat that bounds it. In outline: R1 rests on the
positioning of the separation bar between two measured nulls and,
independently of the geometry, on the rubric-specificity null of
Section~\ref{sec:res-rubric-null}; R2 rests on the leaf-size
distribution and the measured identification cost
(Section~\ref{sec:res-cost}); R3 rests on the fixed-point certificate;
and R4 rests on the acceptance gate being expressed in units of its
own objective's paired standard error
(Section~\ref{sec:res-construction}).

\subsection{Three Things the Construction Does Not Guarantee}\label{sec:arch-limits}

These are stated here, before Chapter~\ref{ch:results}, rather than
collected as limitations afterwards.

\begin{enumerate}
\item \textbf{``Chance-corrected'' means corrected against relabelling,
not against anisotropy.} $J$ compares an observed partition against
random relabelling into the same size profile. A cut taken through a
node's own elongation --- splitting one elongated cluster into two
halves of itself --- is not excluded by that correction, because it
also beats a random relabelling. A subset of accepted splits cannot be
distinguished from such a cut on the evidence
$J$ provides. The counts that quantified this exposure --- how many
accepted splits are affected, and what share of leaf-held queries sit
under at least one such ancestor --- are \textsc{void as to values}
(marker defined in Section~\ref{sec:res-construction})
and are not restated here.

\item \textbf{The size floor is a birth constraint.} It is enforced on
the routed partition at split time. Populations drift as later
iterations re-route queries, so about one leaf in a hundred ends
marginally below it --- 1 of 87 in this artifact, with two further
consistent instances at other settings. Any statement of the form
``every leaf holds at least \texttt{minClusterSize} queries'' is false
as written. This is a property of the mechanism, not a defect.

\item \textbf{The certificate tests period 1 only.} It asks whether one
further application of the operator changes anything. An oscillation at
tolerance scale --- a two-cycle --- would report
\texttt{CERTIFIED: false} indefinitely and would not be distinguished
from divergence. A period-$p$ certificate is not implemented.
\end{enumerate}

None of R1--R4 says that the partition helps the arena. R1 and R2 say
the arena can be run inside these cells and its per-cell estimates
identified; R3 and R4 say the object being run inside is fixed and
described. Whether cell-scoped judging produces better judgments is
links~2 and~3 of Section~\ref{sec:research-questions}, and it is
Chapter~\ref{ch:results}.
